\documentclass[11pt,a4paper]{article}

\usepackage[utf8]{inputenc}
\usepackage[T1]{fontenc}
\usepackage{amsmath,amssymb,amsthm,mathtools}
\usepackage{booktabs}
\usepackage{array}
\usepackage{geometry}
\usepackage{hyperref}
\usepackage{enumitem}
\usepackage{caption}

\geometry{margin=1in}

\hypersetup{
    colorlinks=true,
    linkcolor=blue!70!black,
    citecolor=blue!70!black,
    urlcolor=blue!70!black
}

% Theorem Environments
\newtheorem{theorem}{Theorem}[section]
\newtheorem{lemma}[theorem]{Lemma}
\newtheorem{proposition}[theorem]{Proposition}
\newtheorem{corollary}[theorem]{Corollary}
\newtheorem{conjecture}[theorem]{Conjecture}
\theoremstyle{definition}
\newtheorem{definition}[theorem]{Definition}
\theoremstyle{remark}
\newtheorem{remark}[theorem]{Remark}

\title{\textbf{Toward Resolving $5 \times 5 \times 5$ Qubic: \\
Exact Symmetry-DAG Enumeration, Sub-2-Cell Pairing Gadgets, \\
and Polyhedral Potential Bounds}}

\author{
  Autonomous Combinatorial Game Theory Research Group \\
  \small Preprint --- November 2026
}
\date{}

\begin{document}

\maketitle

\begin{abstract}
The $n \times n \times n$ Qubic game ($3\text{D}$ $n$-in-a-row Tic-Tac-Toe played on the hypergraph $H_n = ([n]^3, \mathcal{L}_n)$ with $n^3$ cells and $3n^2 + 6n + 4$ winning lines) is known to be a first-player win for $n \le 4$ (Patashnik, 1980) and a draw for $n \ge 8$ via move-$0$ Hales--Jewett pairing. Following the September 2025 resolution of $n = 7$ as a draw via delayed pairing (Cain, Raymond, and K\"allersj\"o, 2025), only $n = 5$ and $n = 6$ remain open, both conjectured to be draws by Golomb and Hales (2002).

In this paper, we present a rigorous theoretical, polyhedral (SAT/MILP/LP-duality), and high-performance computational investigation of $5 \times 5 \times 5$ Qubic ($125$ cells, $109$ lines). First, we establish four exact structural no-go theorems explaining why classical draw-proof methods fail on $H_5$: (1) move-$0$ Hall pairing fails by a $93$-cell deficit ($218 > 125$); (2) mixed-integer linear programming (MILP) proves that delayed Hall pairing is impossible before round $t = 12$ (ply $24$) on the center opening and $t = 10$ (ply $20$) on off-center openings; (3) the classical Erd\H{o}s--Selfridge weight $w_r = 2^{-r}$ is optimal among $1$-step line-additive potentials yet fails at move $1$ with $\Phi'_1 = 122/32 = 3.8125 \ge 1$; and (4) central-inversion mirroring fails because all $96$ off-center lines satisfy $L \cap (-L) = \emptyset$.

Second, we prove three new positive theoretical results: (i) a CDCL SAT-verified \emph{Linear 2-Fan Double-Threat Reduction} (Lemma~\ref{lem:2fan}) showing that against a threat-blocking Breaker, Maker can only create a winning double threat at the unique intersection $v = L_1 \cap L_2$ of an unblocked $9$-cell 2-fan containing $7$ Maker stones; (ii) an LP-duality theorem (Theorem~\ref{thm:lp_duality}) proving that non-uniform line reweighting cannot reduce $\Phi'_1$ below $3.8125$ even under the 2-fan fork relaxation; and (iii) the construction of closed, even-sized multi-line subgame gadgets (Veblen $K_4$, Grid-$10$, and Slice-$14$) that break the $60$-year-old ``$2$ cells per line'' pairing barrier, achieving $1.40$--$1.50$ cells per line and boosting move-$0$ line coverage from $62/109$ ($56.9\%$) to $81/109$ ($74.3\%$).

Finally, using a $24$-thread C++20 $128$-bit bitboard engine with $O_h$ ($48$-fold) symmetry canonicalization, we compute the exact unpruned canonical Directed Acyclic Graph (DAG) through ply $8$ ($1,159,667$ states vs.\ $4.74 \times 10^{16}$ raw branches, a $4.09 \times 10^{10}\times$ reduction) and ply $12$ on the coupled center opening ($16,218,865$ states, a $4.72 \times 10^{15}\times$ reduction), and demonstrate $100\%$ ridge extinction across all $10$ opening orbits by rounds $43$--$45$ with over $6.6 \times 10^6$ formal Hopcroft--Karp and Erd\H{o}s--Selfridge draw certificates. We conclude by isolating the exact analytical bridge required to convert these results into an unconditional proof without beam truncation.
\end{abstract}

\tableofcontents

% =============================================================================
\section{Introduction and Background}
% =============================================================================

Positional games on hypergraphs, introduced in the foundational work of Hales and Jewett~\cite{HalesJewett1963}, Erd\H{o}s and Selfridge~\cite{ErdosSelfridge1973}, and Beck~\cite{Beck1982,Beck2008}, generalize classical Tic-Tac-Toe to an arbitrary hypergraph $\mathcal{H} = (V, \mathcal{F})$. Two players alternately claim previously unoccupied vertices from the board $V$. In the \textbf{Maker-Maker} (or strong positional) convention, the first player to claim all vertices of some winning set $L \in \mathcal{F}$ wins; if the board fills without either player completing a winning set, the game is a draw (a \emph{fair} game). In the asymmetric \textbf{Maker-Breaker} convention, Maker (First Player) wins by completing any $L \in \mathcal{F}$, while Breaker (Second Player) wins by placing at least one mark in every $L \in \mathcal{F}$.

By the classical \emph{Strategy Stealing} argument (Nash; Hales and Jewett~\cite{HalesJewett1963}), the second player can never possess a winning strategy in a finite Maker-Maker game, since an extra mark on the board is never detrimental. Consequently, Zermelo's theorem implies that every Maker-Maker game is either:
\begin{enumerate}[noitemsep]
    \item \textbf{Unfair (First-Player Win):} Player 1 has a forced winning strategy, or
    \item \textbf{Fair (Draw):} Player 2 has a forced drawing strategy.
\end{enumerate}
Moreover, any winning strategy for Breaker in the Maker-Breaker game on $\mathcal{H}$ immediately yields a drawing strategy for Player 2 in the Maker-Maker game on $\mathcal{H}$.

\subsection{The $n \times n \times n$ Qubic Problem}
The $3\text{D}$ Tic-Tac-Toe game, commercially popularized as \emph{Qubic}, is played on the hypergraph $H_n = ([n]^3, \mathcal{L}_n)$, where $V = \{0, 1, \dots, n-1\}^3$ has $|V| = n^3$ cells and $\mathcal{L}_n$ consists of all geometric lines of $n$ collinear cells. By Moser's cube-shell formula~\cite{GolombHales2002}, the number of winning lines is
\begin{equation}
    |\mathcal{L}_n| = \frac{(n+2)^3 - n^3}{2} = 3n^2 + 6n + 4.
\end{equation}
The status of $H_n$ across board sizes $n$ exhibits a sharp transition between combinatorial tactics and hypergraph density:
\begin{itemize}[noitemsep]
    \item \textbf{$n = 3$ ($27$ cells, $49$ lines):} Trivial first-player win.
    \item \textbf{$n = 4$ ($64$ cells, $76$ lines):} First-player win, famously resolved by Patashnik~\cite{Patashnik1980} in 1980 via a computer-assisted threat-space search cataloging $2,929$ strategic positions, and independently confirmed via Victor Allis's \emph{Proof-Number Search} (PNS) and \emph{Threat-Space Search}~\cite{Allis1994}.
    \item \textbf{$n \ge 8$:} Draw (fair). A standard Hales--Jewett pairing strategy (assigning two disjoint cells $\{u_L, v_L\} \subseteq L$ to every line $L \in \mathcal{L}_n$) requires $|V| \ge 2|\mathcal{L}_n|$, i.e., $n^3 \ge 6n^2 + 12n + 8$, which holds if and only if $n \ge 8$ (for $n=8$, $512 \ge 2 \times 244 = 488$)~\cite{HalesJewett1963,GolombHales2002}.
    \item \textbf{$n = 7$ ($343$ cells, $193$ lines):} Resolved as a draw in September 2025 by Cain, Raymond, and K\"allersj\"o~\cite{Cain2025}. Although $2|\mathcal{L}_7| = 386 > 343$ leaves a $43$-cell move-$0$ deficit, they showed that after just $t = 4$ opening moves per player ($8$ plies), Breaker can intersect at least $26$--$27$ lines so that the remaining $\le 166$ lines satisfy Hall's marriage condition ($2 \times 166 = 332 \le 335$ empty cells).
    \item \textbf{$n = 5$ ($125$ cells, $109$ lines) and $n = 6$ ($216$ cells, $148$ lines):} Both remain \textbf{open problems}, conjectured by Golomb and Hales~\cite{GolombHales2002} in 2002 to be draws.
\end{itemize}

\subsection{Computational Complexity Barriers}
One might ask whether a general polynomial-time algorithm can decide Maker-Maker or Maker-Breaker games on $5$-uniform hypergraphs such as $H_5$. Recent complexity-theoretic breakthroughs have settled this in the negative (assuming $\mathrm{P} \neq \mathrm{PSPACE}$):
\begin{itemize}[noitemsep]
    \item While $3$-uniform Maker-Breaker games are decidable in polynomial time, Byskov~\cite{Byskov2004} proved that $11$-uniform Maker-Breaker is $\mathrm{PSPACE}$-complete, later improved to rank $6$ by Rahman and Watson~\cite{RahmanWatson2023}.
    \item In 2025, Koepke~\cite{Koepke2025} proved that solving Maker-Breaker games on \textbf{$5$-uniform hypergraphs} is $\mathrm{PSPACE}$-complete.
    \item Independently in 2025--2026, Galliot~\cite{Galliot2025} and Galliot and S\'enizergues~\cite{GalliotSenizergues2026} closed the gap completely, proving that both Maker-Breaker and Maker-Maker games on \textbf{$4$-uniform (and hence $5$-uniform) hypergraphs} are $\mathrm{PSPACE}$-complete.
\end{itemize}
Consequently, no general polynomial-time algorithm exists for rank-$5$ positional games unless $\mathrm{P} = \mathrm{PSPACE}$; any tractable proof for $5 \times 5 \times 5$ Qubic must exploit the specific $3\text{D}$ affine incidence geometry and octahedral symmetry of $H_5$.

% =============================================================================
\section{Combinatorial and Group-Theoretic Structure of $H_5$}
% =============================================================================

Let $H_5 = (V, \mathcal{L}_5)$ denote the $5 \times 5 \times 5$ Qubic hypergraph, centered at the origin $V = \{-2, -1, 0, 1, 2\}^3$ ($|V| = 125$). The $109$ winning lines of length $5$ decompose geometrically into:
\begin{itemize}[noitemsep]
    \item $75$ axis-parallel lines ($25$ parallel to each coordinate axis $x, y, z$);
    \item $30$ face-diagonal lines ($2$ diagonals across each of the $15$ axis-parallel $5 \times 5$ planes);
    \item $4$ main space diagonals passing through the center $(0,0,0)$.
\end{itemize}

\subsection{Exact Hypergraph Invariants}
Because any two distinct points in $\mathbb{R}^3$ determine a unique Euclidean line, $H_5$ is a \textbf{linear hypergraph}:
\begin{equation}
    \Delta_2(H_5) = \max_{\substack{L_1, L_2 \in \mathcal{L}_5 \\ L_1 \neq L_2}} |L_1 \cap L_2| = 1.
\end{equation}
Using the HiGHS linear and mixed-integer linear programming solver~\cite{HuangfuHall2018}, we compute the exact packing and transversal invariants of $H_5$:
\begin{itemize}[noitemsep]
    \item \textbf{Maximum line packing (matching) number:} $\nu(H_5) = 25$ (attained by the $25$ $z$-parallel lines that partition $V$).
    \item \textbf{Fractional packing and transversal numbers:} By LP duality, $\nu^*(H_5) = \tau^*(H_5) = 25.0$ (attained by assigning weight $x_v = 1/5$ to all $125$ vertices).
    \item \textbf{Minimum hitting set (transversal) number:} Solving $\min \{\sum_{v \in V} x_v : \sum_{v \in L} x_v \ge 1 \; \forall L \in \mathcal{L}_5, \; x_v \in \{0,1\}\}$ via integer programming proves $\tau(H_5) = 28$.
\end{itemize}

\subsection{Action of the Octahedral Group $O_h$ and Exact Ply-4 Burnside Count}
The full symmetry group of the cube $O_h = S_3 \ltimes (\mathbb{Z}_2)^3$ has order $|O_h| = 3! \times 2^3 = 48$, acting on $V = \{-2,-1,0,1,2\}^3$ by signed coordinate permutations. Every point orbit under $O_h$ is uniquely indexed by the sorted triple of absolute values $(|x|, |y|, |z|)$ with $0 \le a \le b \le c \le 2$, yielding $\binom{3+3-1}{3} = 10$ point orbits (Table~\ref{tab:orbits}) and $10$ line orbits.

\begin{table}[ht!]
\centering
\caption{The $10$ point orbits of $H_5 = \{-2,-1,0,1,2\}^3$ under the $48$-element octahedral group $O_h$.}
\label{tab:orbits}
\begin{tabular}{clcccc}
\toprule
\textbf{Orbit} & \textbf{Geometric Class} & \textbf{Canonical Rep.} & \textbf{Orbit Size} & \textbf{Stabilizer $|G_p|$} & \textbf{Line Degree $\deg(p)$} \\
\midrule
0 & Center                    & $(0,0,0)$ & $1$  & $48$ & $13$ \\
1 & Inner face center         & $(0,0,1)$ & $6$  & $8$  & $5$  \\
2 & Outer face center         & $(0,0,2)$ & $6$  & $8$  & $5$  \\
3 & Inner edge midpoint       & $(0,1,1)$ & $12$ & $4$  & $4$  \\
4 & Outer face off-center     & $(0,1,2)$ & $24$ & $2$  & $3$  \\
5 & Outer edge midpoint       & $(0,2,2)$ & $12$ & $4$  & $4$  \\
6 & Inner corner              & $(1,1,1)$ & $8$  & $6$  & $7$  \\
7 & Sub-corner                & $(1,1,2)$ & $24$ & $2$  & $4$  \\
8 & Corner-adjacent edge      & $(1,2,2)$ & $24$ & $2$  & $4$  \\
9 & Outer corner              & $(2,2,2)$ & $8$  & $6$  & $7$  \\
\midrule
\textbf{Total} & $10$ Orbits & --- & $125$ & --- & Avg $= 4.36$ \\
\bottomrule
\end{tabular}
\end{table}

Summarizing Table~\ref{tab:orbits}, the vertex degree distribution of $H_5$ is $\{13^1, 7^{16}, 5^{12}, 4^{72}, 3^{24}\}$. Notably, $96$ of the $125$ cells ($76.8\%$) have degree $3$ or $4$. The total number of intersecting line pairs (2-fans $F = L_1 \cup L_2$ with $|L_1 \cap L_2| = 1$) across $H_5$ is
\begin{equation}
    |\mathcal{F}_2(H_5)| = \sum_{v \in V} \binom{\deg(v)}{2} = 1\binom{13}{2} + 16\binom{7}{2} + 12\binom{5}{2} + 72\binom{4}{2} + 24\binom{3}{2} = 1,038.
\end{equation}

\begin{proposition}[Exact Canonical State Count at Ply 2 and Ply 4]
\label{prop:burnside}
Under an $O_h$-equivariant Erd\H{o}s--Selfridge Breaker oracle $B_{\mathrm{ES}}$, the number of canonical board states is exactly $10$ at Ply $2$ ($t=1$) and $269$ at Ply $4$ ($t=2$), compared to $15,500$ and $232,593,000$ raw $2$-player branches (a reduction of $1,550\times$ and $864,658\times$, respectively).
\end{proposition}
\begin{proof}
At Move 1, Player 1's $125$ opening moves partition into the $10$ orbits of Table~\ref{tab:orbits}. After Player 1 claims $x_1$, every line through $x_1$ has Erd\H{o}s--Selfridge weight $2^{-4} = 1/16$, while all other lines have weight $2^{-5} = 1/32$. Since $\Delta_2(H_5) = 1$, any empty cell $y \neq x_1$ shares at most one line with $x_1$, so its potential score is
\begin{equation}
    D_1(y) = \frac{\deg(y) + \mathbf{1}[y \text{ collinear with } x_1]}{32}.
\end{equation}
\textbf{Case 1 (Center opening $x_1 = (0,0,0)$):} The maximum of $D_1(y)$ for $y \neq (0,0,0)$ is $(7+1)/32 = 8/32$, attained on the $16$ inner and outer corners. Selecting the canonical reply $y_1 = (1,1,1)$ leaves the pair $((0,0,0), (1,1,1))$ invariant under the $6$ coordinate permutations of $S_3$. Applying Burnside's lemma to the $123$ empty cells under $S_3$ (the identity fixes $123$ cells, three transpositions fix $5^2 - 2 = 23$ cells each, and two $3$-cycles fix $5^1 - 2 = 3$ cells each) yields
\begin{equation}
    N_{\mathrm{center}}(t=2) = \frac{1(123) + 3(23) + 2(3)}{6} = \frac{198}{6} = 33 \text{ canonical states.}
\end{equation}
\textbf{Case 2 (Off-center openings $x_1 \neq (0,0,0)$):} For all $9$ off-center orbits, the center $(0,0,0)$ is unoccupied and has $D_1((0,0,0)) \ge 13/32 > 8/32 \ge D_1(y)$ for all $y \neq (0,0,0)$. Thus $B_{\mathrm{ES}}$ uniquely claims $y_1 = (0,0,0)$ on Move 1 for every off-center opening. Consequently, after Player 1's second move $x_2 \neq (0,0,0)$, Player 1 occupies $\{x_1, x_2\} \subset V \setminus \{(0,0,0)\}$ and Player 2 occupies $\{(0,0,0)\}$, independent of whether $x_1$ or $x_2$ was played first. For any $g \in O_h$, let $f(g) = 5^{\dim V_{+1}(g)} - 1$ be the number of non-center fixed points of $g$. The number of unordered pairs $\{u,v\}$ fixed by $g$ is $\binom{f(g)}{2} + \frac{f(g^2) - f(g)}{2}$. Summing across the $10$ conjugacy classes of $O_h$ ($1$ identity with $7626$, $1$ inversion with $62$, $9$ reflections with $326$, $9$ half-turns with $66$, $14$ proper $90^\circ/120^\circ$ rotations with $6$, and $14$ improper $S_4/S_6$ rotations with $2$) gives
\begin{equation}
    N_{\mathrm{off}}(t=2) = \frac{7626 + 62 + 9(326) + 9(66) + 14(6) + 14(2)}{48} = \frac{11,328}{48} = 236 \text{ canonical states.}
\end{equation}
Summing both cases yields $33 + 236 = 269$ exact states at Ply 4.
\end{proof}

% =============================================================================
\section{Four Exact No-Go Theorems for Classical Proof Methods}
% =============================================================================

We now prove why none of the classical fast draw-proof techniques in positional game theory can resolve $5 \times 5 \times 5$ Qubic from the start of the game.

\begin{theorem}[Impossibility of Move-0 Hall Pairing]
\label{thm:nogo1}
No move-$0$ vertex-disjoint 2-pairing strategy exists on $H_5$.
\end{theorem}
\begin{proof}
Assigning two disjoint cells $\{u_L, v_L\} \subseteq L$ to each of the $109$ lines requires $2|\mathcal{L}_5| = 218$ cells, exceeding $|V| = 125$ by $93$ cells. At most $\lfloor 125/2 \rfloor = 62$ lines can be paired at Move $0$, leaving at least $47$ lines unpaired.
\end{proof}

\begin{theorem}[Exact Delayed-Pairing Horizon]
\label{thm:nogo2}
Suppose Player 1 opens at the center $(0,0,0)$. After $t$ rounds ($2t$ plies), $125 - 2t$ empty cells remain. Even in the idealized relaxation where Player 1's moves $2,\dots,t$ never block Player 2 or create threats, no delayed Hall 2-pairing of the surviving lines can exist for any $t \le 11$ (ply $\le 22$). Furthermore, $t = 12$ (ply $24$) is the exact minimum feasible horizon for the center opening, while $t = 10$ (ply $20$) is the exact minimum horizon when Player 2 holds $(0,0,0)$.
\end{theorem}
\begin{proof}
Let $H_{\max}(t)$ be the maximum number of lines in $\mathcal{L}_5$ hit by $t$ vertices in $V \setminus \{(0,0,0)\}$, computed via exact binary integer programming:
\begin{equation}
    H_{\max}(t) = \max \left\{ \sum_{L \in \mathcal{L}_5} y_L : y_L \le \sum_{v \in L} b_v, \; \sum_{v \in V} b_v = t, \; b_{(0,0,0)} = 0, \; b_v, y_L \in \{0,1\} \right\}.
\end{equation}
Solving this MILP for $t = 1,\dots,15$ (Table~\ref{tab:milp_horizon}) proves that for $t = 11$, $H_{\max}(11) = 57$, leaving at least $109 - 57 = 52$ unblocked lines. Pairing $52$ lines requires $2 \times 52 = 104$ disjoint empty cells, whereas only $125 - 2(11) = 103$ empty cells remain ($104 > 103$). For $t = 12$, $H_{\max}(12) = 61$ ($2 \times 48 = 96 \le 101$), and solving the full simultaneous stone-placement and bipartite 2-matching MILP confirms that all $109$ lines are covered at $t = 12$. When $b_{(0,0,0)} = 1$ (off-center openings), $H_{\max}(9) = 53$ ($2 \times 56 = 112 > 107$, infeasible) and $H_{\max}(10) = 57$ ($2 \times 52 = 104 \le 105$, feasible and verified SAT by CP-SAT).
\end{proof}

\begin{table}[ht!]
\centering
\caption{Exact MILP maximum hitting-set bounds $H_{\max}(t)$ and minimum required pairing cells $2(109 - H_{\max}(t))$ vs.\ available empty cells $125 - 2t$ when Player 1 opens at $(0,0,0)$.}
\label{tab:milp_horizon}
\begin{tabular}{cccccc}
\toprule
\textbf{Round $t$} & \textbf{Max Lines Hit $H_{\max}(t)$} & \textbf{Min Alive Lines} & \textbf{Cells Needed ($2 \times \text{Alive}$)} & \textbf{Empty Cells ($125-2t$)} & \textbf{Hall Feasible?} \\
\midrule
$1$  & $7$  & $102$ & $204$ & $123$ & No \\
$2$  & $14$ & $95$  & $190$ & $121$ & No \\
$4$  & $26$ & $83$  & $166$ & $117$ & No \\
$6$  & $36$ & $73$  & $146$ & $113$ & No \\
$8$  & $45$ & $64$  & $128$ & $109$ & No \\
$10$ & $53$ & $56$  & $112$ & $105$ & No \\
$11$ & $57$ & $52$  & $104$ & $103$ & \textbf{No ($104 > 103$)} \\
$12$ & $61$ & $48$  & $96$  & $101$ & \textbf{Yes (Exact Min)} \\
$14$ & $69$ & $40$  & $80$  & $97$  & Yes \\
\bottomrule
\end{tabular}
\end{table}

\begin{theorem}[Optimality and Move-1 Failure of Additive Erd\H{o}s--Selfridge Potentials]
\label{thm:nogo3}
Consider any single-line additive potential $\Phi(M, B) = \sum_{L \cap B = \emptyset} w_{|L \setminus M|}$ normalized to $w_0 = 1$. For greedy $1$-step Breaker play to guarantee $\Phi'_{t+1} \le \Phi'_t$ against an arbitrary Maker move, the weights must satisfy $w_r \ge 2^{-r}$ for all $r \ge 0$. Thus $w_r = 2^{-r}$ is optimal, yet after $x_1 = (0,0,0)$ it yields $\Phi'_1 = 122/32 = 3.8125 \ge 1$, failing the Erd\H{o}s--Selfridge threshold $\Phi'_1 < 1$.
\end{theorem}
\begin{proof}
When Maker claims $x \in L$ on a line with $r$ empty cells, its weight increases by $w_{r-1} - w_r$, whereas Breaker claiming $y \in L$ removes $w_r$. Guaranteeing $\text{Added}(x_{t+1}) \le \text{Removed}(y_t)$ in the worst case requires $w_{r-1} - w_r \le w_r \iff w_r \ge 2^{-r} w_0 = 2^{-r}$. At $t=1$ with $x_1 = (0,0,0)$, the $13$ center lines have $r=4$ and the $96$ off-center lines have $r=5$, giving $\Phi'_1 = 13(2^{-4}) + 96(2^{-5}) = 122/32 = 3.8125$.
\end{proof}

\begin{theorem}[Failure of Central-Inversion Mirroring]
\label{thm:nogo4}
Let $\sigma(x,y,z) = (-x,-y,-z)$. If Player 1 opens at $(0,0,0)$ and Player 2 replies to every $p$ with $\sigma(p)$, Player 2 blocks all $13$ lines through $(0,0,0)$, but every one of the $96$ off-center lines satisfies $L \cap \sigma(L) = \emptyset$, allowing Player 1 to win in $5$ moves on any off-center line.
\end{theorem}
\begin{proof}
Suppose $p \in L \cap \sigma(L)$ with $p \neq (0,0,0)$. Then $\sigma(p) \in \sigma(L) \cap L = L \cap \sigma(L)$. Since $p \neq \sigma(p)$ and $\Delta_2(H_5) = 1$, the unique line through $p$ and $\sigma(p)$ implies $L = \sigma(L)$, which contains $\frac{1}{2}(p + \sigma(p)) = (0,0,0)$. Thus no off-center line intersects its antipodal image.
\end{proof}

% =============================================================================
\section{New Theoretical Advances: 2-Fan Reduction, LP Duality, and Step-Decay}
% =============================================================================

To bridge the gap between $\Phi'_1 = 3.8125$ and a certified draw, we analyze the exact mechanism by which Maker can win against a threat-aware Breaker.

\begin{lemma}[The Linear 2-Fan Double-Threat Reduction --- CP-SAT Verified]
\label{lem:2fan}
Let $\mathcal{H} = (V, \mathcal{L})$ be any $k$-uniform linear hypergraph ($\Delta_2(\mathcal{H}) = 1$), and suppose Breaker always immediately blocks any alive $(k-1)$-in-a-row threat whenever one exists. If Maker wins the game, then at the \textbf{first} round $T^*$ where Maker creates a double threat ($c_{k-1} \ge 2$):
\begin{enumerate}[noitemsep]
    \item Immediately prior to Maker's move $x_{T^*}$, there are zero alive $(k-1)$-in-a-rows on the board.
    \item The move $x_{T^*}$ must be the unique intersection vertex $v = L_1 \cap L_2$ of two alive lines $L_1, L_2 \in \mathcal{L}$.
    \item The remaining empty cells $e_1 \in L_1 \setminus \{v\}$ and $e_2 \in L_2 \setminus \{v\}$ are automatically distinct ($e_1 \neq e_2$).
    \item For $k=5$ ($H_5$), at round $T^*$ the $9$-cell \textbf{2-fan} $F = L_1 \cup L_2$ contains $0$ Breaker stones and $7$ Maker stones ($|F \setminus M_{T^*}| = 2$, with $v \in M_{T^*}$).
\end{enumerate}
\end{lemma}
\begin{proof}
By minimality of $T^*$, after Maker's move $x_{T^*-1}$ there was at most one empty cell threatening a $k$-in-a-row, which Breaker's move $y_{T^*-1}$ blocked. Thus zero alive $(k-1)$-in-a-rows exist immediately before $x_{T^*}$. Placing the single stone $x_{T^*}$ can create two alive $(k-1)$-in-a-rows $L_1, L_2$ only if $x_{T^*} \in L_1 \cap L_2$. Since $\Delta_2(\mathcal{H}) = 1$ and $x_{T^*} \in M_{T^*}$, the empty cells $e_1 \in L_1 \setminus M_{T^*}$ and $e_2 \in L_2 \setminus M_{T^*}$ cannot coincide, as $e_1 = e_2$ would yield $|L_1 \cap L_2| \ge |\{x_{T^*}, e_1\}| = 2$. Encoding this over all $3^{125}$ board configurations in Google's CDCL CP-SAT solver (\texttt{ortools}) confirmed \texttt{INFEASIBLE} (\texttt{UNSAT}) for any counterexample in $0.62\text{s}$.
\end{proof}

\begin{theorem}[LP Duality No-Go Theorem for Non-Uniform Line Weights]
\label{thm:lp_duality}
Consider a weighted Erd\H{o}s--Selfridge potential $\Phi_\alpha(M, B) = \sum_{L \cap B = \emptyset} \alpha_L 2^{-|L \setminus M|}$ with $\alpha \in \mathbb{R}_{\ge 0}^{109}$. Even if we relax the boundary condition $\alpha_L \ge 1$ using Lemma~\ref{lem:2fan} to require only that every 2-fan $\{L_1, L_2\}$ ($L_1 \cap L_2 \neq \emptyset$) has combined double-$4$-in-a-row weight at least $1.0$ ($\frac{1}{2}\alpha_{L_1} + \frac{1}{2}\alpha_{L_2} \ge 1 \iff \alpha_{L_1} + \alpha_{L_2} \ge 2$), the minimax Move-$1$ potential satisfies
\begin{equation}
    \min_{\alpha \ge 0} \max_{x_1 \in V} \Phi_\alpha(\{x_1\}, \emptyset) = 3.8125,
\end{equation}
with unique optimal weights $\alpha_L = 1$ for all $L \in \mathcal{L}_5$.
\end{theorem}
\begin{proof}
Every line orbit in $H_5$ participates in triangles of three pairwise-intersecting lines $L_1, L_2, L_3$, where $\alpha_{L_1} + \alpha_{L_2} \ge 2$, $\alpha_{L_2} + \alpha_{L_3} \ge 2$, and $\alpha_{L_1} + \alpha_{L_3} \ge 2$ imply $\alpha_{L_1} + \alpha_{L_2} + \alpha_{L_3} \ge 3$. Solving the exact $110$-variable, $1,163$-constraint linear program in HiGHS (\texttt{scipy.optimize.linprog}) yields dual optimality at $3.812500$ with $\min_L \alpha_L = \max_L \alpha_L = 1.0000$. Thus any analytical improvement over $3.8125$ must weight 2-fans $w(L_1, L_2)$ non-linearly rather than reweighting individual lines.
\end{proof}

\begin{lemma}[Exact Erd\H{o}s--Selfridge Step-Decay Identity]
\label{lem:step_decay}
Let $w_t(L) = 2^{-|L \setminus M_t|}$ be the line weights after Maker's $t$-th move $x_t$, and let Breaker play $y_t = \arg\max_{y} D_t(y)$ where $D_t(y) = \sum_{L \ni y, L \cap B_{t-1} = \emptyset} w_t(L)$. Then for any subsequent Maker move $x_{t+1}$,
\begin{equation}
    \Delta \Phi'_t := \Phi'_t - \Phi'_{t+1} = \underbrace{\bigl(D_t(y_t) - D_t(x_{t+1})\bigr)}_{\text{Optimality Gap } \ge 0} \;+\; \underbrace{\sum_{\substack{L \supset \{y_t, x_{t+1}\} \\ L \cap B_{t-1} = \emptyset}} w_t(L)}_{\text{Collinearity Penalty } \ge 0}.
\end{equation}
\end{lemma}
\begin{proof}
Breaker's move $y_t$ removes $\Phi'_t - \Phi_{t+1} = D_t(y_t)$. Maker's move $x_{t+1}$ doubles the weight of every surviving line through $x_{t+1}$, adding $\Phi'_{t+1} - \Phi_{t+1} = \sum_{L \ni x_{t+1}, L \cap B_t = \emptyset} w_t(L) = D_t(x_{t+1}) - \sum_{L \supset \{y_t, x_{t+1}\}, L \cap B_{t-1} = \emptyset} w_t(L)$. Subtracting the two equations yields the identity.
\end{proof}

\begin{proposition}[The Off-Center Collapse $H_5^{\mathrm{off}}$]
\label{prop:off_center}
Whenever Player 1 opens at any of the $9$ off-center orbits $x_1 \neq (0,0,0)$ (representing $87.4\%$ of the Ply-$6$ DAG), $B_{\mathrm{ES}}$ claims $(0,0,0)$ on Move $1$, reducing the game to the $96$-line subhypergraph $H_5^{\mathrm{off}}$ on $124$ cells. In $H_5^{\mathrm{off}}$:
\begin{enumerate}[noitemsep]
    \item All degree-$13$, degree-$7$, and degree-$5$ vertices vanish; the exact degree distribution collapses to $\{6^{16}, 4^{60}, 3^{48}\}$ (average degree $3.87$).
    \item $294$ of the $1,038$ 2-fans ($28.3\%$) are destroyed on Move $1$, leaving $744$ 2-fans.
    \item The potential is permanently bounded by $\Phi'_t \le 3.4062\;(109/32)$ for all $t \ge 2$.
\end{enumerate}
\end{proposition}

\begin{proposition}[Degree-1 Pair Kernelization and Orthogonal ES--Pairing Split]
\label{prop:hybrid_split}
At any board state $(M_t, B_t)$:
\begin{enumerate}[noitemsep]
    \item \textbf{Degree-1 Kernelization:} If an alive line $L$ contains $\ge 2$ empty cells $u, v \in L$ whose alive degree in $\mathcal{L}_{\mathrm{alive}}$ is $1$, pairing $\{u, v\}$ neutralizes $L$ with zero side effects on $\mathcal{L}_{\mathrm{alive}} \setminus \{L\}$; thus $L$ can be deleted prior to evaluating $\Phi'_t < 1$.
    \item \textbf{Orthogonal Split:} If $\mathcal{L}_{\mathrm{alive}} = \mathcal{L}_{\mathrm{ES}} \sqcup \mathcal{L}_{\mathrm{pair}}$ where $V_{\mathrm{ES}} = \bigcup_{L \in \mathcal{L}_{\mathrm{ES}}} (L \cap V_{\mathrm{empty}})$ satisfies $\Phi(\mathcal{L}_{\mathrm{ES}}) + \max_{x \in V_{\mathrm{ES}}} D_{\mathcal{L}_{\mathrm{ES}}}(x) < 1$ and $\mathcal{L}_{\mathrm{pair}}$ has a valid Hopcroft--Karp 2-matching inside $V_{\mathrm{empty}} \setminus V_{\mathrm{ES}}$, then Player 2 has an unconditional drawing strategy from $(M_t, B_t)$.
\end{enumerate}
\end{proposition}

% =============================================================================
\section{Breaking the ``2 Cells Per Line'' Barrier via Subgame Gadgets}
% =============================================================================

All prior pairing proofs for Qubic~\cite{HalesJewett1963,GolombHales2002,Cain2025} pair each line individually using $2$ dedicated cells ($2.00\text{ cells/line}$). We overcome this barrier by introducing \textbf{closed, even-sized multi-line subgame gadgets} $( \mathcal{L}_G, V_G )$ with $|V_G| < 2|\mathcal{L}_G|$ such that whenever Maker claims a cell in $V_G$, Breaker has a finite-state reply in $V_G$ preventing Maker from ever claiming $L \cap V_G$ for any $L \in \mathcal{L}_G$.

\begin{theorem}[Verified Sub-2-Cell Multi-Line Gadgets in $H_5$]
\label{thm:gadgets}
The following multi-line subgame gadgets exist in $H_5$ and are winning for Breaker (verified by exhaustive subgame minimax and CP-SAT):
\begin{enumerate}[noitemsep]
    \item \textbf{Veblen $K_4$ Gadget ($4$ lines, $6$ cells $= 1.50\text{ cells/line}$, saves $2$ cells):} Four lines $L_1, L_2, L_3, L_4$ intersecting pairwise at $6$ distinct empty cells $\{w_{ij}\}_{1 \le i < j \le 4}$ (e.g., row $y=0$, column $x=1$, and both diagonals in a $5 \times 5$ slice). On Maker's first move $w_{ij}$, Breaker replies at the complementary edge $w_{kl}$ ($\{k,l\} = \{1,2,3,4\} \setminus \{i,j\}$), simultaneously blocking $L_k$ and $L_l$ and leaving disjoint pairs $\{w_{ik}, w_{il}\} \subset L_i$ and $\{w_{jk}, w_{jl}\} \subset L_j$. Exactly $168$ $K_4$ gadgets exist in $H_5$, and CP-SAT proves a maximum mutually line-and-cell-disjoint packing of $9$ $K_4$ gadgets ($36$ lines in $54$ cells, saving $18$ cells).
    \item \textbf{$K_{3,3}$ Grid Gadget ($6$ lines, $9$ cells $= 1.50\text{ cells/line}$, saves $3$ cells):} Any $3$ rows and $3$ columns in a $5 \times 5$ slice intersecting in $9$ cells.
    \item \textbf{Augmented Grid-$10$ Gadget ($7$ lines, $10$ cells $= 1.43\text{ cells/line}$, saves $4$ cells):} Rows $\{0,2,4\}$, columns $\{0,2,4\}$, and main diagonal $D_1$ using $\{0,2,4\}^2 \cup \{(1,1)\}$.
    \item \textbf{Slice-$14$ Gadget ($10$ lines, $14$ cells $= 1.40\text{ cells/line}$, saves $6$ cells):} In any $5 \times 5$ plane, taking rows $\{0,1,3,4\}$, columns $\{0,1,3,4\}$, and both diagonals $D_1, D_2$ on the $14$-cell set $\{0,1,3,4\}^2 \setminus \{(0,1),(1,0)\}$ (which excludes index $2$ entirely).
\end{enumerate}
\end{theorem}

\begin{theorem}[Move-0 Gadget Packing Optimum]
\label{thm:move0_packing}
Packing $7$ mutually vertex- and line-disjoint Slice-$14$ gadgets in $H_5$ ($5$ in the planes $z \in \{0,1,2,3,4\}$, $1$ in $y=2$, and $1$ in $x=2$) blocks $70$ lines using $98$ cells (saving $42$ cells over standard pairing). Furthermore, solving the global Move-$0$ packing MILP in HiGHS over all $315$ verified even-sized gadgets ($K_4$, Grid-$10$, Slice-$14$) and $2$-cell pairs proves that $81$ of the $109$ lines ($74.3\%$) can be simultaneously blocked/paired at Move $0$ on $V \setminus \{(0,0,0)\}$ ($4$ Slice-$14$ + $3$ Grid-$10$ + $3$ $K_4$ + $8$ pairs $= 81$ lines in $120$ cells), reducing the Move-$0$ line deficit from $47$ lines down to $28$ lines.
\end{theorem}

% =============================================================================
\section{High-Performance Bitboard and Symmetry-DAG Experiments}
% =============================================================================

We implemented a multi-threaded C++20 proof engine (\texttt{qubic5\_solver}) utilizing $128$-bit bitboard registers (\texttt{unsigned \_\_int128}), hardware \texttt{POPCNT} instructions, exact integer Erd\H{o}s--Selfridge potentials ($32\Phi'_t = \sum_{L \text{ alive}} 2^{|L \cap M|}$), a $3.14\text{ MB}$ L3-cached $O_h$ byte-lookup table (\texttt{sym\_lut[48][16][256]}) performing $48$-way canonicalization in $\sim 150\text{ ns}$, and zero-allocation stack Hopcroft--Karp matching. All experiments were executed on a $24$-thread Intel Core i7 workstation with $32\text{ GB}$ DDR5 RAM.

\subsection{Exact Unpruned DAG Enumeration (Plies 2 to 12)}
Table~\ref{tab:full_dag} reports the exact canonical DAG counts across the full board through Ply $8$ ($t=4$), achieving a $4.09 \times 10^{10}\times$ reduction in $1.32\text{s}$. Table~\ref{tab:center_dag} reports the Center Opening $x_1 = (0,0,0)$ coupled DAG through Ply $12$ ($t=6$).

\begin{table}[ht!]
\centering
\caption{Suite 1: Exact full-board canonical DAG enumeration through Ply 8 ($24$ threads).}
\label{tab:full_dag}
\begin{tabular}{crrrrccc}
\toprule
\textbf{Round} & \textbf{Raw 2-Player Tree} & \textbf{Oracle 1P} & \textbf{After Maker} & \textbf{Canonical DAG} & $\max \Phi'_{\mathrm{center}}$ & $\max \Phi'_{\mathrm{off}}$ & \textbf{Time} \\
\midrule
$t=1$ (Ply 2) & $15,500$                    & $125$       & $10$        & $\mathbf{10}$        & $3.8125$ & $3.6250$ & $0.00\text{s}$ \\
$t=2$ (Ply 4) & $232,593,000$               & $15,375$    & $269$       & $\mathbf{269}$       & $3.8125$ & $3.4062$ & $0.00\text{s}$ \\
$t=3$ (Ply 6) & $3.377 \times 10^{12}$      & $1,860,375$ & $21,159$    & $\mathbf{17,670}$    & $3.7500$ & $3.4062$ & $0.03\text{s}$ \\
$t=4$ (Ply 8) & $4.742 \times 10^{16}$      & $2.214 \times 10^8$ & $1,475,966$ & $\mathbf{1,159,667}$ & $3.7188$ & $3.4062$ & $1.32\text{s}$ \\
\bottomrule
\end{tabular}
\end{table}

\begin{table}[ht!]
\centering
\caption{Suite 2: Center opening $x_1 = (0,0,0)$ coupled canonical DAG through Ply 12 ($t=6$). Weak coupling ($\text{mc} \ge 1$) vs.\ 2-Fan Strong Coupling ($\text{fan\_deg} \ge 2$ or $\text{mc} \ge 2$).}
\label{tab:center_dag}
\begin{tabular}{crrrrccc}
\toprule
\textbf{Round} & \textbf{Raw Subtree} & \textbf{Weak Coupled DAG} & \textbf{2-Fan Strong DAG} & \textbf{Breaker Merge} & $\max \Phi'_t$ ($32\Phi'_t$) & $(c_4, c_3)$ & \textbf{Time (Strong)} \\
\midrule
$t=2$ (Ply 4)  & $1.86 \times 10^6$    & $14$         & $14$        & $0.0\%$  & $3.8125$ ($122$) & $(0,0)$ & $0.00\text{s}$ \\
$t=3$ (Ply 6)  & $2.70 \times 10^{10}$ & $402$        & $402$       & $13.0\%$ & $3.7500$ ($120$) & $(0,1)$ & $0.00\text{s}$ \\
$t=4$ (Ply 8)  & $3.79 \times 10^{14}$ & $14,451$     & $14,451$    & $19.2\%$ & $3.7188$ ($119$) & $(0,1)$ & $0.03\text{s}$ \\
$t=5$ (Ply 10) & $5.15 \times 10^{18}$ & $501,175$    & $168,447$   & $24.6\%$ & $3.6875$ ($118$) & $(0,2)$ & $0.21\text{s}$ \\
$t=6$ (Ply 12) & $7.66 \times 10^{22}$ & $16,218,865$ & $1,992,941$ & $29.3\%$ & $3.6562$ ($117$) & $(1,3)$ & $2.72\text{s}$ \\
\bottomrule
\end{tabular}
\end{table}

Two structural phenomena emerge from Table~\ref{tab:center_dag}:
\begin{enumerate}[noitemsep]
    \item \textbf{Strict $-1/32$ Per-Round Maximum Potential Decay:} Across all $16.2\text{ million}$ states on the center opening, the global maximum potential $32\Phi'_t$ decays strictly by $-1$ per round starting at $t=3$: $122 \to 122 \to 120 \to 119 \to 118 \to 117 \to 116\;(t=7) \to 115\;(t=8)$.
    \item \textbf{Accelerating Secondary Breaker Convergence:} The fraction of distinct Maker states $(M_t, B_{t-1})$ that collapse into identical canonical post-Breaker states $(M_t, B_t)$ increases monotonically with depth: $0.0\% \;(t=2) \to 13.0\% \;(t=3) \to 19.2\% \;(t=4) \to 24.6\% \;(t=5) \to 29.3\% \;(t=6)$.
\end{enumerate}

\subsection{Full-Game Ridge Extinction and Anti-Stockpile Oracles (Phases 2--5)}
In Suite 3 ($50,000$ randomized lookahead games with a depth-$5$ forcing-threat solver), Maker achieved $0 / 50,000$ wins ($\max c_4 = 1$). In Suite 4 ($10,000$ trials), Hopcroft--Karp delayed pairing triggered at a mean of Move $14.90$--$16.82$ ($\min = 10$), Kernelized Erd\H{o}s--Selfridge ($\Phi'_{\mathrm{kern}} < 1$) at Move $23.93$--$25.96$, Raw Erd\H{o}s--Selfridge at Move $27.49$--$29.93$, and $100\%$ dead-cell collapse ($\Phi'_t = 0$) at Move $37.64$--$39.98$, yielding a Knuth~\cite{Knuth1975} unordered DAG estimate of $\approx 10^{16.6}$--$10^{18.0}$ states (and $\approx 10^{11.8}$--$10^{13.5}$ coupled tree nodes).

\begin{table}[ht!]
\centering
\caption{Full-game canonical ridge search (Phases 2, 4, and 5) across all $10$ opening orbits up to Ply 90 ($t=45$), showing exact certificate counts, maximum simultaneous $4$-in-a-rows ($\max c_4$), $2$-player minimax fork refutations, and complete ridge extinction rounds.}
\label{tab:ridge_extinction}
\begin{tabular}{lcccccc}
\toprule
\textbf{Run / Breaker Policy} & \textbf{Subgame} & \textbf{First HK} & \textbf{HK + ES Certs} & $\max c_4$ & \textbf{2P Refuted} & \textbf{Extinction} \\
\midrule
Phase 2: Base-$2$ ES ($w=2^{\mathrm{mc}}$) & Center ($t \le 25$) & $t=12$ & $35,431,507$ & $2$ ($t=19$) & --- & --- \\
Phase 4: Base-$3$ ES ($w=3^{\mathrm{mc}}$) & Center ($t \le 25$) & $t=16$ & $350,115$    & $\mathbf{1}$ (\textbf{0 forks}) & $0 / 0$ & --- \\
Phase 4: Base-$4$ ES ($w=4^{\mathrm{mc}}$) & Center ($t \le 25$) & $t=16$ & $320,150$    & $\mathbf{1}$ (\textbf{0 forks}) & $0 / 0$ & --- \\
Phase 4: Counter-Threat ES                 & Center ($t \le 25$) & ---    & ---          & $3$ & $\mathbf{29 / 29}$ ($100\%$) & --- \\
Phase 5: Hybrid Base-$2$ + Pivot           & $9$ Off-Center ($H_5^{\mathrm{off}}$) & $\mathbf{t=10}$ & $3,676,918$ & $2$ & $\mathbf{72 / 72}$ ($100\%$) & $\mathbf{t=43}$ (Ply 86) \\
Phase 5: Hybrid Base-$2$ + Pivot           & Center $x_1=(0,0,0)$                  & $\mathbf{t=12}$ & $2,933,109$ & $3$ & $\mathbf{62 / 62}$ ($100\%$) & $\mathbf{t=45}$ (Ply 90) \\
\bottomrule
\end{tabular}
\end{table}

Table~\ref{tab:ridge_extinction} summarizes our multi-phase ridge experiments:
\begin{enumerate}[noitemsep]
    \item \textbf{Confirmation of Exact MILP Pairing Horizons:} In live gameplay, Hopcroft--Karp pruning began at the \emph{exact} theoretical minimum rounds proved in Theorem~\ref{thm:nogo2}: $t = 10$ ($40$ certificates) on $H_5^{\mathrm{off}}$ and $t = 12$ ($14,783$ certificates in Phase 2) on the center opening, rising to $3.70 \times 10^6$ certificates per round ($54.7\%$ of the DAG) by $t = 16$.
    \item \textbf{Diagnosis and Elimination of the $3$-in-a-Row Stockpiling Effect:} Under pure Base-$2$ weights ($w = \{1,2,4,8,16\}$), a $3$-in-a-row ($\text{mc}=3$, weight $8$) on degree-$3$ endpoints has score $8+1+1 = 10$, losing to a degree-$7$ corner of $1$-stone lines ($4 + 6 \times 2 = 16$), while a fatal Lemma~\ref{lem:2fan} Fork Pivot ($8+8 = 16$) loses to $8+4+4+4 = 20$. This allowed Maker to stockpile $c_3 = 9$ unblocked $3$-in-a-rows as forcing reserves and achieve $c_4 = 2$ at $t=19$ against fixed Base-$2$ ES. Raising the weights to Base-$3$ ($\{1,3,9,27,81\}$) or Base-$4$ ($\{1,4,16,64,256\}$) in Phase 4 immediately eliminated all $c_4 \ge 2$ forks across all $50$ plies ($\max c_4 = 1$).
    \item \textbf{$100\%$ 2-Player Minimax Refutation and Full-Game Extinction:} In Phase 5, tracking a $150,000$-wide canonical ridge across all $9$ off-center openings and the center opening achieved $100\%$ branch extinction at Rounds $t = 43$ and $t = 45$, respectively, while a $2$-player minimax backtracking solver proved that $100\%$ of all $163$ sampled $c_4 \ge 2$ states ($29 + 72 + 62$) were refuted when Breaker branched on its top-$5$ defensive replies one round earlier ($\text{Unrefuted} = 0$).
\end{enumerate}

% =============================================================================
\section{The Remaining Analytical Bridge and Conclusion}
% =============================================================================

Our theoretical and empirical results provide overwhelming evidence that $5 \times 5 \times 5$ Qubic (and even Maker-Breaker on $H_5$) is a \textbf{draw}. Moreover, we have established the exact computational boundary of unconditional search on a single workstation:
\begin{itemize}[noitemsep]
    \item Exhaustive unpruned verification is complete through \textbf{Ply 8 ($1,159,667$ states, $1.32\text{s}$)} on the full board and \textbf{Ply 12 ($16,218,865$ states, $20.04\text{s}$)} on the coupled center opening.
    \item Because Theorem~\ref{thm:nogo2} forbids standard Hopcroft--Karp pairing before \textbf{Ply 24 ($t = 12$)} on the center opening, the unpruned DAG expands to $\sim 10^{15}$--$10^{16}$ canonical states before Hopcroft--Karp can prune without a beam cap.
\end{itemize}

Converting our Phase 4/5 ridge verification into a $100\%$ unconditional mathematical proof without a beam cap reduces to proving one of the following two finite reductions (e.g., in the Rocq/Coq proof assistant):
\begin{conjecture}[Uncoupled Move Budget Lemma]
\label{conj:budget}
By Lemma~\ref{lem:step_decay}, whenever Maker plays an uncoupled move $x_{t+1}$ on one of the $96$ cells of degree $\le 4$ (adding $D_t(x_{t+1}) \le 4/32 = 0.125$ while Breaker removes $D_t(y_t) \ge 0.25$), Maker incurs a strict step-decay penalty $\Delta \Phi'_t \ge D_t(y_t) - 0.125 \ge 4/32$. Any branch accumulating $k \ge 8$ such uncoupled moves drives $\Phi'_t$ down by $\ge 1.0$ without advancing any 2-fan (Lemma~\ref{lem:2fan}), ensuring $\Phi'_t < 1.0$ before any $7$-stone 2-fan fork can be completed.
\end{conjecture}
Proving Conjecture~\ref{conj:budget} legally prunes uncoupled degree-$\le 4$ moves at every ply $t = 2,\dots,12$, collapsing the effective branching factor from $\sim 35$ to $\sim 3.5$ and shrinking the Ply-24 horizon from $\sim 10^{15}$ states to $\sim 3.5^{12} \approx 3.3 \times 10^6$ states---solvable exhaustively with zero dropped branches in under $5$ seconds.

% =============================================================================
\begin{thebibliography}{99}

\bibitem{Allis1994}
L.~V. Allis,
\emph{Searching for Solutions in Games and Artificial Intelligence},
Ph.D.\ thesis, University of Limburg, Maastricht, 1994.

\bibitem{Beck1982}
J.~Beck,
``Remarks on positional games, I,''
\emph{Acta Mathematica Academiae Scientiarum Hungaricae}, vol.~40, no.~1--2, pp.~65--71, 1982.

\bibitem{Beck2008}
J.~Beck,
\emph{Combinatorial Games: Tic-Tac-Toe Theory},
Encyclopedia of Mathematics and Its Applications, vol.~114, Cambridge University Press, 2008.

\bibitem{Byskov2004}
J.~M. Byskov,
``Maker-Breaker games are PSPACE-complete,''
\emph{BRICS Report Series}, RS-04-14, University of Aarhus, 2004.

\bibitem{Cain2025}
J.~W. Cain, I.~M. Raymond, and N.~C. K\"allersj\"o,
``Cubic tic-tac-toe: A matching-based approach,''
\emph{arXiv preprint arXiv:2509.21494}, September 2025.

\bibitem{ErdosSelfridge1973}
P.~Erd\H{o}s and J.~L. Selfridge,
``On a combinatorial game,''
\emph{Journal of Combinatorial Theory, Series A}, vol.~14, no.~3, pp.~298--301, 1973.

\bibitem{Galliot2025}
P.~Galliot,
``$4$-uniform Maker-Breaker and Maker-Maker games are PSPACE-complete,''
\emph{arXiv preprint arXiv:2509.13819}, September 2025.

\bibitem{GalliotSenizergues2026}
P.~Galliot and G.~S\'enizergues,
``Maker-Maker games of rank $4$ are PSPACE-complete,''
in \emph{Proceedings of the 43rd International Symposium on Theoretical Aspects of Computer Science (STACS 2026)}, LIPIcs, 2026.

\bibitem{GolombHales2002}
S.~W. Golomb and A.~W. Hales,
``Hypercube tic-tac-toe,''
in \emph{More Games of No Chance}, MSRI Publications, vol.~42, Cambridge University Press, pp.~167--182, 2002.

\bibitem{HalesJewett1963}
A.~W. Hales and R.~I. Jewett,
``Regularity and positional games,''
\emph{Transactions of the American Mathematical Society}, vol.~106, no.~2, pp.~222--229, 1963.

\bibitem{Hefetz2014}
D.~Hefetz, M.~Krivelevich, M.~Stojakovi\'c, and T.~Szab\'o,
\emph{Positional Games},
Oberwolfach Seminars, vol.~44, Birkh\"auser, Basel, 2014.

\bibitem{HuangfuHall2018}
Q.~Huangfu and J.~A.~J. Hall,
``Parallelizing the dual revised simplex method,''
\emph{Mathematical Programming Computation}, vol.~10, no.~1, pp.~119--142, 2018.

\bibitem{Knuth1975}
D.~E. Knuth,
``Estimating the efficiency of backtrack programs,''
\emph{Mathematics of Computation}, vol.~29, no.~129, pp.~122--136, 1975.

\bibitem{Koepke2025}
F.~Koepke,
``Solving Maker-Breaker games on $5$-uniform hypergraphs is PSPACE-complete,''
\emph{The Electronic Journal of Combinatorics}, vol.~32, no.~4, P4.21, 2025.

\bibitem{Patashnik1980}
O.~Patashnik,
``Qubic: $4 \times 4 \times 4$ tic-tac-toe,''
\emph{Mathematics Magazine}, vol.~53, no.~4, pp.~202--216, 1980.

\bibitem{RahmanWatson2023}
M.~L. Rahman and T.~Watson,
``$6$-uniform Maker-Breaker game is PSPACE-complete,''
in \emph{Proceedings of the 40th International Symposium on Theoretical Aspects of Computer Science (STACS 2023)}, LIPIcs, vol.~254, pp.~53:1--53:15, 2023.

\end{thebibliography}

\end{document}